
\section{The next interface to quantum integration}
\label{nr16-next}

The classical background and its response are now controlled
on a uniform homogeneous small-field domain. The minimized
action has a true nonlinear local response and a coarse
Hessian close to the flat Maxwell precision. Adding an
auxiliary positive coarse penalty gives a controlled inverse,
but that penalty is not a physical mass and this inverse
is not yet the covariance produced by the complete quantum
renormalization step.

Exponential response alone does not give strict background
support for polymer activities. The physical Hessian,
gauge-coordinate factors and multiple-insertion local
decomposition require the separate constructions described
below. Regions outside the chart and preservation of uniform
complex-source bounds under quantum iteration remain essential.

Notebook 17 now proves a relative energy comparison for the actual
first-step reduced Hessian, retaining constraint curvature, and
joint Gaussian bad-block estimates on the fixed-step domain. Its
homogeneous flat current-scale theorem covers all previous averaging
depths. Notebook 18 handles overlapping local interaction supports
and strict background locality under a small local precision
decomposition. Notebook 19 supplies the actual physical
pivot Hessian and local measure factors. Notebook 20 closes
the bare Wilson small-field integration. Notebook 21 supplies
strict support for the minimized classical action, with
arbitrary fixed cardinality reserve, and extends one
classical-reference integration uniformly through earlier
refinement. The next-coarse support transport, inherited
quantum interactions and regional domains remain separate.

The continuum obligations remain those of Notebook 01:
a common nontrivial gauge-invariant observable hierarchy,
Euclidean covariance and positivity in the limit, appropriate
normalization and Yang--Mills identities, and the additional
vacuum and mass-gap properties. The classical theorems here
remove specific prerequisites; they do not establish those
quantum conclusions.

\begin{thebibliography}{9}
\bibitem{prior16}
\emph{Toward a four-dimensional Yang--Mills construction},
research Notebooks 8--15 in the present workspace,
16 October 2026. The weighted flat extension and inverse,
uniform nonlinear constraint chart, Wilson force, constrained
minimum and ordinary-average transfer are the prerequisites.
\bibitem{source16}
E.~Sussman and collaborators,
\href{https://github.com/ethanwsussman/Yang--Mills-LLM-experiment}{Yang--Mills-LLM-experiment},
deposited raw accounts I--III at commit
\texttt{e392ef45013a581664bab76b90867170b30fe10e}.
The provenance limitations of Notebook 15 remain in force.
\end{thebibliography}
\end{document}
